\section{Second-Degree Equations and Conic Sections}

Lines arise from equations in which the coordinates occur only to the first
power. The next natural class consists of equations in which terms of total
degree two are allowed.

\begin{definitionbox}
A \emph{second-degree equation} in the Cartesian plane has the form
\[
Ax^2+Bxy+Cy^2+Dx+Ey+F=0,
\]
where $A,B,C,D,E,F\in\mathbb{R}$ and at least one of $A$, $B$, or $C$ is
nonzero.
\]
Its solution set is the collection of all points $(x,y)\in\mathbb{R}^2$ whose
coordinates satisfy the equation.
\end{definitionbox}

The terms
\[
Ax^2+Bxy+Cy^2
\]
form the quadratic part. The remaining terms translate the curve and affect its
position, while the quadratic part largely determines its geometric type.

\subsection{A First Classification}

For the equation
\[
Ax^2+Bxy+Cy^2+Dx+Ey+F=0,
\]
define
\[
\Delta=B^2-4AC.
\]

\begin{theorembox}
For a non-degenerate real conic, the sign of $\Delta$ gives the following
classification:
\[
\begin{array}{c|c}
\Delta<0 & \text{ellipse type, including circles},\\
\Delta=0 & \text{parabola type},\\
\Delta>0 & \text{hyperbola type}.
\end{array}
\]
\end{theorembox}

This classification is preserved by rotations and translations of the
coordinate system. A rotation can remove the mixed term $Bxy$, while a
translation can move the centre or vertex to a simpler position.

\begin{remarkbox}
The sign of $\Delta$ identifies the conic type only after degenerate cases are
separated. A second-degree equation may also describe a single point, no real
points, one line, or two lines. These possibilities are discussed below.
\end{remarkbox}

\subsection{Circle}

A circle with centre $(h,k)$ and radius $r>0$ is described by
\[
(x-h)^2+(y-k)^2=r^2.
\]
After expansion,
\[
x^2+y^2-2hx-2ky+h^2+k^2-r^2=0.
\]
Here $A=C=1$ and $B=0$, so
\[
\Delta=0^2-4\cdot1\cdot1=-4<0.
\]
A circle is therefore a special ellipse-type conic in which the coefficients of
$x^2$ and $y^2$ are equal and there is no mixed term after a suitable rotation.

\begin{center}
\begin{tikzpicture}[scale=0.82]
    \draw[axis] (-3.8,0) -- (4.2,0) node[right] {$x$};
    \draw[axis] (0,-3.3) -- (0,3.7) node[above] {$y$};
    \coordinate (C) at (1,0.6);
    \draw[subset] (C) circle (2);
    \fill[geometricSubsetColor] (C) circle (2pt) node[below right] {$(h,k)$};
    \draw[helper] (C) -- ++(35:2) node[midway,above] {$r$};
\end{tikzpicture}
\end{center}

\subsection{Ellipse}

The standard ellipse centred at the origin has equation
\[
\frac{x^2}{a^2}+\frac{y^2}{b^2}=1,
\qquad a>0,\quad b>0.
\]
Equivalently,
\[
b^2x^2+a^2y^2-a^2b^2=0.
\]
Thus $A=b^2$, $B=0$, and $C=a^2$, so
\[
\Delta=-4a^2b^2<0.
\]
The parameters $a$ and $b$ are the semi-axis lengths. When $a=b$, the ellipse
becomes a circle.

\begin{center}
\begin{tikzpicture}[scale=0.85]
    \draw[axis] (-4.2,0) -- (4.5,0) node[right] {$x$};
    \draw[axis] (0,-3.0) -- (0,3.3) node[above] {$y$};
    \draw[subset] (0,0) ellipse (3.1 and 1.8);
    \draw[helper] (0,0) -- (3.1,0) node[midway,below] {$a$};
    \draw[helper] (0,0) -- (0,1.8) node[midway,left] {$b$};
\end{tikzpicture}
\end{center}

\subsection{Parabola}

A standard vertical parabola has equation
\[
y=ax^2,
\qquad a\neq0,
\]
or equivalently
\[
ax^2-y=0.
\]
Here $A=a$, $B=0$, and $C=0$, hence
\[
\Delta=0.
\]
The sign of $a$ determines whether the parabola opens upward or downward. A
translated form is
\[
y-k=a(x-h)^2,
\]
whose vertex is $(h,k)$.

\begin{center}
\begin{tikzpicture}[scale=0.82]
    \draw[axis] (-3.8,0) -- (3.9,0) node[right] {$x$};
    \draw[axis] (0,-1.0) -- (0,5.2) node[above] {$y$};
    \draw[subset,domain=-2.2:2.2,samples=100] plot (\x,{0.85*\x*\x});
    \fill[geometricSubsetColor] (0,0) circle (2pt) node[below right] {vertex};
\end{tikzpicture}
\end{center}

\subsection{Hyperbola}

A standard hyperbola centred at the origin has equation
\[
\frac{x^2}{a^2}-\frac{y^2}{b^2}=1,
\qquad a>0,\quad b>0.
\]
Equivalently,
\[
b^2x^2-a^2y^2-a^2b^2=0.
\]
Here $A=b^2$, $B=0$, and $C=-a^2$, so
\[
\Delta=4a^2b^2>0.
\]
Its asymptotes are
\[
y=\frac{b}{a}x,
\qquad
y=-\frac{b}{a}x.
\]

\begin{center}
\begin{tikzpicture}[scale=0.78]
    \draw[axis] (-5.0,0) -- (5.2,0) node[right] {$x$};
    \draw[axis] (0,-4.0) -- (0,4.2) node[above] {$y$};
    \draw[helper] (-4.6,-3.07) -- (4.6,3.07);
    \draw[helper] (-4.6,3.07) -- (4.6,-3.07);
    \draw[subset,domain=1.55:4.5,samples=100] plot (\x,{1.1*sqrt(\x*\x/2.25-1)});
    \draw[subset,domain=1.55:4.5,samples=100] plot (\x,{-1.1*sqrt(\x*\x/2.25-1)});
    \draw[subset,domain=-4.5:-1.55,samples=100] plot (\x,{1.1*sqrt(\x*\x/2.25-1)});
    \draw[subset,domain=-4.5:-1.55,samples=100] plot (\x,{-1.1*sqrt(\x*\x/2.25-1)});
\end{tikzpicture}
\end{center}

\subsection{The Mixed Term and Rotation}

The term $Bxy$ indicates that the natural axes of the conic may be rotated
relative to the chosen Cartesian axes. For example,
\[
x^2+2xy+y^2=1
\]
can be written as
\[
(x+y)^2=1,
\]
so it is not a non-degenerate ellipse: it describes the two parallel lines
\[
x+y=1,
\qquad
x+y=-1.
\]

For a genuine rotated conic, one introduces new coordinates $(u,v)$ obtained by
a rotation through an angle $\theta$:
\[
x=u\cos\theta-v\sin\theta,
\qquad
y=u\sin\theta+v\cos\theta.
\]
Choosing $\theta$ appropriately removes the $uv$ term. When $A\neq C$, one may
choose $\theta$ so that
\[
\tan(2\theta)=\frac{B}{A-C}.
\]
The resulting equation can then be classified in the new coordinates using the
standard forms above.

\subsection{Degenerate Second-Degree Equations}

Not every second-degree equation describes a non-degenerate conic. Factorisation
or completing squares may reveal a simpler set.

\begin{examplebox}
The equation
\[
x^2-y^2=0
\]
factors as
\[
(x-y)(x+y)=0.
\]
Its solution set is the union of two intersecting lines:
\[
y=x
\qquad\text{or}\qquad
y=-x.
\]
\end{examplebox}

\begin{center}
\begin{tikzpicture}[scale=0.75]
    \draw[axis] (-4.0,0) -- (4.2,0) node[right] {$x$};
    \draw[axis] (0,-3.8) -- (0,4.0) node[above] {$y$};
    \draw[subset] (-3.5,-3.5) -- (3.5,3.5);
    \draw[subset] (-3.5,3.5) -- (3.5,-3.5);
\end{tikzpicture}
\end{center}

Other common degenerate cases include
\[
x^2+y^2=0,
\]
which gives only the point $(0,0)$,
\[
x^2+y^2+1=0,
\]
which has no real solutions, and
\[
x^2=0,
\]
which describes the single line $x=0$ counted algebraically twice.

\begin{summarybox}
Second-degree equations have the general form
\[
Ax^2+Bxy+Cy^2+Dx+Ey+F=0.
\]
For non-degenerate real conics, the discriminant
\[
\Delta=B^2-4AC
\]
separates the three principal types:
\[
\Delta<0\text{ gives ellipse type},\qquad
\Delta=0\text{ gives parabola type},\qquad
\Delta>0\text{ gives hyperbola type}.
\]
Translations and rotations reduce many equations to standard forms, while
factorisation and completing squares expose degenerate cases.
\end{summarybox}
